\subsection{赋范代数中元素的谱}

在本小节中，我们假定 K = C。

定理 1. --- 设 A 是一个幺 Banach 代数，并设 \(x \in \mathrm{A}\)。

\begin{enumerate}
\def\labelenumi{\alph{enumi})}
\item
  集合 \(\mathrm{Sp}_{\mathrm{A}}(x)\) 是 C 的一个紧子集；
\item
  谱半径 \(\varrho(x)\) 是 \(\mathbf{C}\) 中以 0 为中心且包含 \(\mathrm{Sp}_{\mathrm{A}}(x)\) 的最小闭圆盘的半径；
\item
  预解式 \(\lambda \mapsto \mathrm{R}(x,\lambda) = (\lambda -x)^{-1}\)（\(x\) 的预解式）在 \(\mathbf{C} - \mathrm{Sp}_{\mathrm{A}}(x)\) 中全纯且在无穷远处为 0。而且，对每个整数 \(k\geqslant 0\)，我们有
下列公式：

\[
\frac {\partial^ {k}}{\partial \lambda^ {k}} \mathrm{R} (x, \lambda) = (- 1) ^ {k} k! \mathrm{R} (x, \lambda) ^ {k + 1};
\]

\item
  对每个复数 \(\lambda\)（满足 \(|\lambda| > 1 / \varrho(x)\)），有

\[
\mathrm{R} (x, \lambda) = \sum_ {n = 0} ^ {+ \infty} \lambda^ {- n - 1} x ^ {n}.
\]

\end{enumerate}

谱 \(\mathrm{Sp}_{\mathrm{A}}(x)\) 的补集是 A 的可逆元所成的群 G 在从 C 到 A 的连续映射 \(\lambda \mapsto x - \lambda\) 下的原像；由 I, p.~23 的命题 3, b)，谱 \(\mathrm{Sp}_{\mathrm{A}}(x)\) 是 C 的一个闭子集。此外，设 \(\lambda \in C\) 满足 \(|\lambda| > \varrho(x)\)。我们有 \(\lambda - x = \lambda(1 - \lambda^{-1}x)\)。由于 \(\varrho(\lambda^{-1}x) = |\lambda|^{-1}\varrho(x) < 1\)，元素 \(1 - \lambda^{-1}x\)，从而元素 \(\lambda - x\)，可逆，且

\[
\mathrm{R} (x, \lambda) = (\lambda - x) ^ {- 1} = \sum_ {n = 0} ^ {\infty} \lambda^ {- n - 1} x ^ {n}\tag{5}
\]

(I, p.~22, 命题 2)。特别地，\(\lambda \notin \mathrm{Sp}_{\mathrm{A}}(x)\)。这就证明 \(\mathrm{Sp}_{\mathrm{A}}(x)\) 包含在以 0 为中心、\(\varrho(x)\) 为半径的圆盘中。因而 \(\mathrm{Sp}_{\mathrm{A}}(x)\) 是紧的。公式 (5) 还表明，\(x\) 的预解式在闭圆盘 \(\Delta_{\varrho(x)}\)（以 0 为中心、\(\varrho(x)\) 为半径）的补集中有定义且全纯，并在无穷远处趋于 0。

设 \(\lambda_0\in \mathbf{C} - \mathrm{Sp}_{\mathrm{A}}(x)\)。令 \(y = \lambda_0 - x\)。设 \(\mu \in \mathbf{C}\) 满足 \(|\lambda_0 - \mu | < \| y^{-1}\|^{-1}\)。我们有

\[
\mu - x = y - (\lambda_ {0} - \mu) = y (1 - (\lambda_ {0} - \mu) y ^ {- 1}),
\]

因而 \(\mu - x\) 可逆，且其逆为

\[
(\mu - x) ^ {- 1} = y ^ {- 1} \sum_ {n = 0} ^ {\infty} (\lambda_ {0} - \mu) ^ {n} y ^ {- n}\tag{6}
\]

(由 I, p.~22 的命题 2)。于是 \(x\) 的预解式在以 \(\lambda_0\) 为中心、\(\| y^{-1} \|^{-1}\) 为半径的开圆盘中有定义且全纯。因而 \(x\) 的预解式是从 \(\mathbf{C} - \mathrm{Sp}_{\mathrm{A}}(x)\) 到 A 中的一个全纯映射。

I, p.~4 的公式 (1) 蕴涵 \(\frac{\partial}{\partial\lambda}\mathrm{R}(x,\lambda) = -\mathrm{R}(x,\lambda)^2\)，由此，对 \(k\) 用归纳法，

\[
\frac {\partial^ {k}}{\partial \lambda^ {k}} \mathrm{R} (x, \lambda) = (- 1) ^ {k} k! \mathrm{R} (x, \lambda) ^ {k + 1}.
\]

设 \(a > 0\) 是使 \(\mathrm{Sp}_{\mathrm{A}}(x)\) 包含在闭圆盘 \(\Delta_{a}\)（以 0 为中心、\(a\) 为半径）中的实数。函数 \(\lambda \mapsto (\lambda^{-1} - x)^{-1}\) 于是在 \(0 < |\lambda| < a^{-1}\) 上有定义且全纯，且当 \(\lambda\) 趋于 0 时趋于 0。在这个开圆盘（以 0 为中心、\(a^{-1}\) 为半径）上延拓此全纯函数的唯一连续函数是全纯的 (VAR, R1, 3.3.9)，因而定义它的级数 (5) 的收敛半径 \(\geqslant a^{-1}\) (VAR, R1, 3.2.9)。由 I, p.~22 的命题 2，我们因此有 \(a \geqslant \varrho(x)\)。

注. --- 1) 幺 Banach 代数中元素的谱可以是 C 的任意一个非空紧子集 F (参见 I, p. 17 的例 3；对 A = C(F; C) 与 f ∈ A（F 到 C 中的典范包含），有 \(\mathrm{Sp}_{\mathrm{A}}(f) = F\))。

\begin{enumerate}
\def\labelenumi{\arabic{enumi})}
\setcounter{enumi}{1}
\tightlist
\item
  设 A 是一个幺 Banach 代数，并设 \(x \in A\)。由 I, p.~24 的定理 1，\(\mathbf{C} - \mathrm{Sp}_{\mathrm{A}}(x)\) 是 C 的一个开子集，故是局部连通的。于是 \(\mathbf{C} - \mathrm{Sp}_{\mathrm{A}}(x)\) 的连通分支都是开的。由定理 1，这些连通分支之一包含由 \(\lambda \in C\)（满足 \(|\lambda| > \varrho(x)\)）所成的集合；所有其他的连通分支因此都是有界的。
\end{enumerate}

推论 1. --- 设 A 是一个非零的赋范幺代数。对每个 \(x \in \mathrm{A}\)，谱 \(\mathrm{Sp}_{\mathrm{A}}(x)\) 非空。

先假定 A 完备。若有 \(\mathrm{Sp}(x)=\varnothing\)，则 x 的预解式就会在 C 中全纯且在无穷远处为 0，从而恒等于零 (VAR, R., 3.3.6, p.~29)。由于 \(\mathrm{R}(x,\lambda)=(\lambda-x)^{-1}\) 可逆，就会得出 1=0，从而 \(A=\{0\}\)。

在一般情形，设 \(\widehat{\mathrm{A}}\) 是 A 的完备化代数；关系 \(\mathrm{Sp}_{\mathrm{A}}(x) = \varnothing\) 将蕴涵 \(\mathrm{Sp}_{\widehat{\mathrm{A}}} (x) = \varnothing\)，从而 \(\widehat{\mathrm{A}} = \{0\}\)，故 \(\mathrm{A} = \{0\}\)。

推论 2 (Gelfand–Mazur 定理). --- 设 A 是 C 上的一个赋范代数。若 A 是一个域，则 A = C · 1。

若 \(x \in A\)，则存在 \(\lambda \in \mathbf{C}\) 使 \(x - \lambda\) 不可逆 (推论 1)，因而 \(x - \lambda = 0\)，即 \(x \in \mathbf{C} \cdot 1\)。

推论 3. --- 设 A 是一个幺 Banach 代数，x 是 A 的一个可逆元素，满足 \(\|x\|=\|x^{-1}\|=1\)。则 \(\mathrm{Sp}(x)\subset\mathbf{U}\)。

设 \(\Delta\) 是 \(\mathbf{C}\) 中以 0 为中心、1 为半径的圆盘。由定理 1 b) 以及 \(\varrho(x) \leqslant \|x\|\) 这一事实，有 \(\operatorname{Sp}(x) \subset \Delta\)。同样，\(\operatorname{Sp}(x)^{-1} = \operatorname{Sp}(x^{-1}) \subset \Delta\)，由此得推论 (参见 I, p.~2, 注 4)。

推论 4. --- 设 E 是一个复 Banach 空间，\(\mathcal{L}(\mathrm{E})\) 是 E 的连续自同态所成的 Banach 代数，A 是 \(\mathcal{L}(\mathrm{E})\) 的一个非零子代数，使得 E 是一个单 A-伪模 (A, II, p.~176, 附录)。

\begin{enumerate}
\def\labelenumi{\alph{enumi})}
\item
  设 u 是 E 的一个自同态，不必连续，它与 A 交换。则 u 是一个位似。
\item
  设 \(u\) 是 \(\mathbf{E}\) 的一个自同态，不必连续。对每个整数 \(n \geqslant 1\) 及每个 \((\xi_1, \ldots, \xi_n) \in \mathrm{E}^n\)，存在 \(v \in A\)，使得

\[
(v (\xi_ {1}), \dots , v (\xi_ {n})) = (u (\xi_ {1}), \dots , u (\xi_ {n})).
\]

\end{enumerate}

我们来证明 \(a\) )。设 \(\widetilde{\mathrm{A}}\) 是由 A 添加单位元而得到的代数。由于 E 是一个单 A-伪模，它是一个单 \(\widetilde{\mathrm{A}}\)-模。

设 B 是 \(\tilde{\mathsf{A}}\) 在 C-向量空间 E 的自同态环中的交换子。代数 B 包含 1，且是 \(\tilde{\mathsf{A}}\)-模 E 的自同态代数。由于 E 是一个单 \(\tilde{\mathsf{A}}\)-模，Schur 引理 (A, VIII, p.~43, 推论) 表明 B 是一个域。

设 \(\xi_0\in \mathrm{E}\) 使 \(\mathrm{A}\xi_0\neq\{0\}\)。于是我们有 \(\mathrm{A}\xi_0 = E\)。对每个 \(u\in \mathrm{B}\)，用 \(\mathrm{A}_u\) 表示由 \(v\in \mathrm{A}\)（使 \(v(\xi_0) = u(\xi_0)\)）所成的集合。这个集合非空，因为 \(\mathrm{A}\xi_0 = E\)。我们令

\[
\| u \| _ {\mathrm{B}} = \inf _ {v \in \mathrm{A} _ {u}} \| v \|.
\]

映射 \(u \mapsto \|u\|_{\mathrm{B}}\) 是 B 上的一个半范数。

我们来证明这个映射是一个范数。设 u 是 B 的一个非零元素。对每个 \(v \in A_{u}\)，有 \(\|v\| \geqslant \|v(\xi_{0})\|/\|\xi_{0}\| = \|u(\xi_{0})\|/\|\xi_{0}\|\)，从而 \(\|u\|_{B} \geqslant \|u(\xi_{0})\|/\|\xi_{0}\|\)。于是只须证明 \(u(\xi_{0}) \neq 0\)。设 \(\xi_{1} \in E\) 使 \(u(\xi_{1}) \neq 0\)。由于 \(A\xi_{0} = E\)，存在 \(w \in A\) 使 \(\xi_{1} = w(\xi_{0})\)。于是 \(wu(\xi_{0}) = uw(\xi_{0}) = u(\xi_{1}) \neq 0\)，因而 \(u(\xi_{0}) \neq 0\)。

另一方面，设 \(u\) 与 \(u'\) 是 B 的元素。对每个 \(\varepsilon > 0\)，存在 \(v, v' \in A\)，使得 \(v(\xi_0) = u(\xi_0), v'(\xi_0) = u'(\xi_0)\)，且 \(\|v\| \leqslant \|u\|_{\mathrm{B}} + \varepsilon\)、\(\|v'\| \leqslant \|u'\|_{\mathrm{B}} + \varepsilon\)。于是我们有 \(vv'(\xi_0) = vu'(\xi_0) = u'v(\xi_0) = u'u(\xi_0)\)，从而

\[
\| u ^ {\prime} u \| _ {\mathrm{B}} \leqslant \| v ^ {\prime} v \| \leqslant \| v \| \| v ^ {\prime} \| \leqslant (\| u \| _ {\mathrm{B}} + \varepsilon) (\| u ^ {\prime} \| _ {\mathrm{B}} + \varepsilon),
\]

最后有 \(\|u'u\|_{B} \leqslant \|u\|_{B} \|u'\|_{B}\)。这表明 B 带上范数 \(u \mapsto \|u\|_{B}\)，是一个赋范代数。由于它是一个域，推论 2 蕴涵 \(B = C \cdot 1\)，这就是所要的结论。

我们证明 \(b\) )。由 \(a\) )，A 在 \(\mathrm{End}_{\mathbf{C}}(\mathrm{E})\) 中的交换子化为 E 的位似。于是其双交换子是 \(\mathrm{End}_{\mathbf{C}}(\mathrm{E})\)。断言 \(b\) ) 于是由 Jacobson 稠密定理 (A, VIII, p.~434 的定理 1) 得出。

推论 5. --- 设 A 是一个 Banach 代数，\(x \in A\)。

\begin{enumerate}
\def\labelenumi{\alph{enumi})}
\item
  \(\mathrm{Sp}^{\prime}(x)\) 是 C 的一个紧子集；
\item
  谱半径 \(\varrho(x)\) 是 C 中以 0 为中心且包含 \(\mathrm{Sp}'(x)\) 的最小闭圆盘的半径；
\item
  要使 \(x\) 是拟幂零的，必须且只需 \(\mathrm{Sp}'(x) = \{0\}\)。
\end{enumerate}

断言 \(a)\) 与 \(b)\) 由定理 1 通过考虑由 A 添加单位元而得到的 Banach 代数得出。断言 \(c)\) 由 \(b)\) 得出。

